\documentclass[10pt, xcolor=table]{beamer}
\usepackage[utf8]{inputenc}
\usepackage[vietnamese]{babel} 
\usepackage{booktabs}          
\usepackage{amsmath}           
\usepackage{amsfonts}          
\usepackage{array}
\usepackage{tabularx}

% Cấu hình Theme slide
\usetheme{Madrid}
\usecolortheme{whale}

% Định nghĩa cột tự động xuống hàng và căn lề trái mượt mà
\newcolumntype{L}[1]{>{\raggedright\arraybackslash}p{#1}}

% Thông tin tiêu đề bài báo cáo
\title[Sparse PGD \& Phòng Thủ]{Báo Cáo Phân Tích: Thuật Toán Tấn công Sparse PGD \& Các Kỹ Thuật Phòng Thủ}
\subtitle{Xử lý ảnh số \& Học máy an toàn}
\author[Môn học: IMP302]{Môn học: IMP302}
\institute[ĐH]{Nguồn dữ liệu dựa trên tài liệu báo cáo phân tích}
\date{\today}

\begin{document}

% Slide 1: Trang tiêu đề
\begin{frame}
    \titlepage
\end{frame}

% ===============================
% PHẦN 1: CƠ SỞ LÝ THUYẾT
% ===============================
\section{Cơ sở lý thuyết}

% Slide 1.1: Thuật toán Cốt lõi
\begin{frame}{1.1 Cơ sở Lý thuyết: Thuật toán Proposed-TopkPGD}
    \begin{block}{Tấn công Sparse PGD dựa trên Gradient đề xuất}
        Thuật toán Sparse PGD giới hạn số lượng điểm ảnh thay đổi trong phạm vi tỷ lệ $k$ (k-ratio) trên tổng số điểm ảnh. \\
        \textbf{Các bước tối ưu chính:}
        \begin{itemize}
            \item \textbf{Tính Gradient:} $g_t = \text{grad}_{x_t} L(f(x_t), y)$
            \item \textbf{Lọc Top-k:} $M_t(c, i, j) = 1$ nếu $|g_t| \ge \tau$ (Chỉ kích hoạt các tọa độ pixel có đạo hàm lớn nhất, đóng băng các pixel còn lại).
            \item \textbf{Cập nhật nhiễu:} $x'_{t+1} = x_t + \alpha \cdot \text{sign}(g_t) \odot M_t$
            \item \textbf{Lập lịch $k$ động:} $k_t = \max(1, \text{int}(k_{\max} \cdot (1 - t/T)))$ giúp thưa hóa dần dần lượng pixel nhiễu theo thời gian.
        \end{itemize}
    \end{block}
\end{frame}

% Slide 1.2: Phương pháp phòng thủ xử lý ảnh
\begin{frame}{1.2 Cơ sở Lý thuyết: Các Phương pháp Phòng thủ Tiền xử lý}
    Mục tiêu của tiền xử lý đầu vào (Input Preprocessing) là phá vỡ cấu trúc không gian của nhiễu đối kháng thưa trước khi đưa ảnh vào mạng CNN:
    \vspace{0.2cm}
    \begin{itemize}
        \item \textbf{Median Smoothing (Lọc trung vị):} Triệt tiêu nhiễu thưa cục bộ bằng cách thay thế pixel bằng giá trị trung vị của các pixel lân cận. Rất hiệu quả với nhiễu muối tiêu (salt-and-pepper) hay Sparse Attack.
        \item \textbf{Bit Depth Reduction (Giảm bit màu):} Lượng hóa kênh màu (ví dụ từ 8-bit xuống 3-bit), làm mất đi phần biên độ tinh vi của nhiễu đối kháng.
        \item \textbf{JPEG Compression (Nén JPEG):} Loại bỏ các thành phần nhiễu ở miền tần số cao (high-frequency) thông qua biến đổi DCT.
        \item \textbf{Random Noise (Nhiễu ngẫu nhiên):} Cộng thêm nhiễu Gaussian ngẫu nhiên để làm lu mờ cấu trúc có định hướng của kẻ tấn công.
    \end{itemize}
\end{frame}

% Slide 1.3: Tổng hợp Cơ chế Phòng thủ (Định nghĩa)
\begin{frame}{1.3 Cơ sở Lý thuyết: Tổng hợp các Cơ chế Phòng thủ Bền vững}
    \begin{itemize}
        \item \textbf{Adversarial Training (Huấn luyện đối kháng):} Huấn luyện mô hình trực tiếp trên các mẫu đối kháng (như PGD) được sinh ra liên tục trong từng vòng lặp. Đây là cách trực tiếp nhất để tăng cường Robustness.
        \item \textbf{TRADES:} Một biến thể của Adversarial Training, sử dụng KL Divergence để tối ưu hóa sự cân bằng (trade-off) giữa độ chính xác trên ảnh sạch (Clean Accuracy) và Robustness.
        \item \textbf{Randomized Smoothing (Làm mịn ngẫu nhiên):} Cộng nhiễu Gaussian nhiều lần vào ảnh và lấy dự đoán chiếm đa số (Majority vote). Cung cấp độ bền bỉ có chứng nhận toán học (certified robustness).
        \item \textbf{Feature Denoising (Khử nhiễu đặc trưng):} Thay vì tiền xử lý ảnh gốc, phương pháp này nhúng các khối lọc nhiễu (Non-local, Median filter) vào tận sâu bên trong cấu trúc CNN để lọc sạch activation maps.
    \end{itemize}
\end{frame}

% ===============================
% PHẦN 2: THỰC NGHIỆM TẤN CÔNG
% ===============================
\section{Thực nghiệm Tấn công}

% Slide 2.1: Kết quả thử nghiệm trên Mô hình Chuẩn
\begin{frame}{2.1 Thực nghiệm Tấn công cơ bản: Mô hình Chuẩn}
    \begin{block}{Độ chính xác trên ảnh sạch (\textit{Clean Accuracy})}
        \centering \textbf{94.80\%}
    \end{block}
    \centering
    \resizebox{\textwidth}{!}{%
    \begin{tabular}{lccccc}
        \toprule
        \textbf{Phương thức tấn công} & \textbf{k-ratio} & \textbf{Accuracy} & \textbf{ASR} & \textbf{Sparsity} & \textbf{PSNR (dB)} \\
        \midrule
        Clean (Ảnh sạch) & -   & 94.80\% & 0.00\%   & 100.0\% & Vô cùng ($\infty$) \\
        FGSM             & -   & 30.50\% & 67.99\%  & 0.3\%   & 30.14 \\
        BIM              & -   & 0.00\%  & 100.00\% & 2.6\%   & 33.22 \\
        PGD              & -   & 0.00\%  & 100.00\% & 0.1\%   & 32.95 \\
        \midrule
        Sparse PGD       & 0.1 & 8.20\%  & 91.47\%  & 60.9\%  & 42.41 \\
        Sparse PGD       & 0.2 & 1.90\%  & 98.03\%  & 37.3\%  & 39.66 \\
        Sparse PGD       & 0.5 & 0.20\%  & 99.79\%  & 8.9\%   & 36.28 \\
        \bottomrule
    \end{tabular}%
    }
\end{frame}

% Slide 2.2: Kết quả thử nghiệm trên Mô hình Đối kháng
\begin{frame}{2.2 Thực nghiệm Tấn công cơ bản: Mô hình Đối kháng (AT)}
    \begin{block}{Độ chính xác trên ảnh sạch (\textit{Clean Accuracy})}
        \centering \textbf{85.90\%}
    \end{block}
    \centering
    \resizebox{\textwidth}{!}{%
    \begin{tabular}{lccccc}
        \toprule
        \textbf{Phương thức tấn công} & \textbf{k-ratio} & \textbf{Accuracy} & \textbf{ASR} & \textbf{Sparsity} & \textbf{PSNR (dB)} \\
        \midrule
        Clean (Ảnh sạch) & -   & 85.90\% & 0.00\%   & 100.0\% & Vô cùng ($\infty$) \\
        FGSM             & -   & 63.00\% & 26.69\%  & 0.2\%   & 30.11 \\
        BIM              & -   & 59.00\% & 31.38\%  & 0.2\%   & 30.43 \\
        PGD              & -   & 59.30\% & 31.02\%  & 0.2\%   & 30.51 \\
        \midrule
        Sparse PGD       & 0.1 & 73.10\% & 14.91\%  & 76.3\%  & 41.15 \\
        Sparse PGD       & 0.2 & 68.20\% & 20.60\%  & 57.9\%  & 38.14 \\
        Sparse PGD       & 0.5 & 62.60\% & 27.17\%  & 19.7\%  & 34.19 \\
        \bottomrule
    \end{tabular}%
    }
\end{frame}

% Slide 2.3: Phân tích số bước lặp PGD
\begin{frame}{2.3 Phân tích quá trình hội tụ của PGD truyền thống}
    \centering
    \small \textbf{PGD ITERATIONS ANALYSIS (Standard vs Robust Models)}
    \vspace{0.1cm}
    
    \resizebox{\textwidth}{!}{%
    \begin{tabular}{lccccc}
        \toprule
        \textbf{Mô hình} & \textbf{Số bước lặp (T)} & \textbf{Accuracy} & \textbf{ASR} & \textbf{Sparsity} & \textbf{PSNR (dB)} \\
        \midrule
        \textbf{Standard} & 1  & 30.30\% & 68.04\% & 0.09\% & 34.51 \\
                          & 3  & 2.40\%  & 97.47\% & 0.11\% & 33.94 \\
                          & 5  & 0.20\%  & 99.79\% & 0.12\% & 33.52 \\
                          & 10 & 0.00\%  & 100.00\% & 0.14\% & 32.95 \\
        \midrule
        \textbf{Robust}   & 1  & 78.40\% & 6.89\%  & 0.05\% & 34.49 \\
                          & 3  & 65.10\% & 22.68\% & 0.11\% & 32.94 \\
                          & 5  & 54.60\% & 35.15\% & 0.17\% & 31.63 \\
                          & 10 & 47.50\% & 43.59\% & 0.20\% & 30.61 \\
        \bottomrule
    \end{tabular}%
    }
    \vspace{0.2cm}
    \begin{itemize}
        \item \textbf{Tốc độ hội tụ của PGD:} Trên mô hình Standard, PGD đạt độ hội tụ rất nhanh chỉ sau 5 bước lặp (ASR đạt $99.79\%$).
        \item \textbf{Độ thưa của nhiễu PGD:} Độ thưa cực kỳ thấp ($<0.2\%$), cho thấy PGD lan truyền nhiễu trên toàn bộ các điểm ảnh (dày đặc).
    \end{itemize}
\end{frame}

% Slide 2.4: Quá trình hội tụ theo số bước lặp của Proposed-TopkPGD
\begin{frame}{2.4 Quá trình hội tụ của Proposed-TopkPGD}
    \centering
    \small \textbf{CONVERGENCE STUDY OF PROPOSED-TOPKPGD}
    \vspace{0.1cm}
    \resizebox{\textwidth}{!}{%
    \begin{tabular}{ccrccrcc}
        \toprule
        & & \multicolumn{3}{c}{\textbf{Mô hình Chuẩn (Standard)}} & \multicolumn{3}{c}{\textbf{Mô hình Đối kháng (Robust)}} \\
        \cmidrule(lr){3-5} \cmidrule(lr){6-8}
        \textbf{k-ratio} & \textbf{T} & \textbf{ASR} & \textbf{Sparsity} & \textbf{PSNR} & \textbf{ASR} & \textbf{Sparsity} & \textbf{PSNR} \\
        \midrule
        & 1  & 30.91\% & 80.12\% & 49.14 & 4.16\%  & 80.08\% & 49.12 \\
        & 3  & 77.32\% & 76.02\% & 43.65 & 15.20\% & 78.42\% & 40.24 \\
        \textbf{0.2} & 5  & 90.51\% & 74.76\% & 41.22 & 22.57\% & 79.38\% & 37.69 \\
        & 10 & 96.94\% & 76.34\% & 39.43 & 22.80\% & 79.48\% & 37.44 \\
        \midrule
        & 1  & 46.20\% & 50.11\% & 45.14 & 5.82\%  & 50.02\% & 45.13 \\
        & 3  & 92.30\% & 44.56\% & 39.52 & 24.11\% & 46.79\% & 36.26 \\
        \textbf{0.5} & 5  & 98.31\% & 40.89\% & 37.20 & 35.63\% & 48.41\% & 33.77 \\
        & 10 & 100.00\%& 42.58\% & 35.52 & 36.94\% & 48.93\% & 33.45 \\
        \bottomrule
    \end{tabular}%
    }
    \vspace{0.1cm}
    \begin{itemize}
        \item Đòn tấn công đạt độ hội tụ ASR rất nhanh chóng chỉ trong vòng \textbf{5 - 7 bước lặp} trên cả hai mô hình. Sau bước 5, ASR thay đổi không đáng kể.
        \item Nhờ cơ chế lập lịch $k$ động, \textbf{độ thưa (Sparsity)} luôn được kiểm soát chặt chẽ ổn định trong suốt quá trình tối ưu.
    \end{itemize}
\end{frame}

% Slide 2.5: Đánh giá SparseFool
\begin{frame}{2.5 Đánh giá nhược điểm thuật toán SOTA SparseFool}
    \begin{block}{Giới thiệu thuật toán SparseFool}
        SparseFool là phương thức tấn công thưa hình học, cố gắng tìm hướng di chuyển tối thiểu đến ranh giới quyết định.
    \end{block}
    \centering
    \resizebox{\textwidth}{!}{%
    \begin{tabular}{lccccccccc}
        \toprule
        \textbf{Mô hình} & \textbf{Accuracy} & \textbf{ASR} & $\mathbf{L_0}$ \textbf{(pixels)} & \textbf{Sparsity} & $\mathbf{L_2}$ & $\mathbf{L_\infty}$ & \textbf{SSIM} & \textbf{PSNR (dB)} & \textbf{Time (s)} \\
        \midrule
        Standard & 57.50\% & 39.35\% & 11.87 & 98.84\% & 1.76 & 0.75 & 0.96 & 29.12 & 241.78 \\
        Robust   & 64.60\% & 23.28\% & 8.93  & 99.13\% & 1.38 & 0.58 & 0.98 & 30.72 & 162.97 \\
        \bottomrule
    \end{tabular}%
    }
    \vspace{0.2cm}
    \begin{itemize}
        \item \textbf{Điểm mạnh - Độ thưa vượt trội:} SparseFool chỉ sửa đổi trung bình khoảng \textbf{9-12 điểm ảnh} (độ thưa đạt mức $>98.8\%$).
        \item \textbf{Điểm yếu - Chi phí tính toán cực lớn:} Thời gian sinh mẫu đối kháng rất lớn (160s - 240s cho 1000 mẫu), khiến SparseFool \textbf{quá chậm} so với các biến thể PGD.
    \end{itemize}
\end{frame}

% Slide 2.6: So sánh chi tiết trên Mô hình Chuẩn
\begin{frame}{2.6 So sánh các Thuật toán Tấn công Thưa (Mô hình Chuẩn)}
    \centering
    \small \textbf{COMPARING SPARSE ATTACKS ON STANDARD MODEL (Iteration = 10)}
    \vspace{0.1cm}
    \resizebox{\textwidth}{!}{%
    \begin{tabular}{clcccc}
        \toprule
        \textbf{k-ratio} & \textbf{Thuật toán} & \textbf{Robust Accuracy} & \textbf{ASR} & \textbf{Sparsity} & \textbf{PSNR (dB)} \\
        \midrule
        & Sparse-PGD & 4.90\% & 94.83\% & 66.11\% & 38.04 \\
        \textbf{0.1} & GreedyFool & 4.80\% & 94.94\% & 64.31\% & 40.20 \\
        & \cellcolor{yellow!20}\textbf{Proposed-TopkPGD} & \cellcolor{yellow!20}\textbf{15.50\%} & \cellcolor{yellow!20}\textbf{83.65\%} & \cellcolor{yellow!20}\textbf{88.23\%} & \cellcolor{yellow!20}\textbf{42.66} \\
        \midrule
        & Sparse-PGD & 1.30\% & 98.63\% & 50.09\% & 36.88 \\
        \textbf{0.2} & GreedyFool & 1.00\% & 98.95\% & 42.65\% & 37.60 \\
        & \cellcolor{yellow!20}\textbf{Proposed-TopkPGD} & \cellcolor{yellow!20}\textbf{2.90\%} & \cellcolor{yellow!20}\textbf{96.94\%} & \cellcolor{yellow!20}\textbf{76.34\%} & \cellcolor{yellow!20}\textbf{39.43} \\
        \midrule
        & Sparse-PGD & 0.30\% & 99.68\% & 30.62\% & 35.31 \\
        \textbf{0.3} & GreedyFool & 0.20\% & 99.79\% & 28.99\% & 36.22 \\
        & \cellcolor{yellow!20}\textbf{Proposed-TopkPGD} & \cellcolor{yellow!20}\textbf{1.00\%} & \cellcolor{yellow!20}\textbf{98.95\%} & \cellcolor{yellow!20}\textbf{64.47\%} & \cellcolor{yellow!20}\textbf{37.62} \\
        \bottomrule
    \end{tabular}%
    }
    \vspace{0.1cm}
    \begin{itemize}
        \item \textbf{Chất lượng ảnh đối kháng xuất sắc:} \textbf{Proposed-TopkPGD} luôn duy trì độ thưa (\textbf{Sparsity}) và độ tương đồng chất lượng ảnh (\textbf{PSNR}) cao hơn đáng kể so với GreedyFool và Sparse-PGD.
        \item \textbf{Cân bằng Hoàn hảo:} Với mức giới hạn gắt gao $k=0.1$, thuật toán đề xuất vẫn giữ được độ thưa lên tới \textbf{88.23\%} (so với $\sim 64\%$ của SOTA) nhưng vẫn bảo đảm sát thương ASR cao (83.65\%).
    \end{itemize}
\end{frame}

% Slide 2.7: So sánh chi tiết trên Mô hình Đối kháng
\begin{frame}{2.7 So sánh các Thuật toán Tấn công Thưa (Mô hình Đối kháng)}
    \centering
    \small \textbf{COMPARING SPARSE ATTACKS ON ROBUST MODEL (Iteration = 10)}
    \vspace{0.1cm}
    \resizebox{\textwidth}{!}{%
    \begin{tabular}{clcccc}
        \toprule
        \textbf{k-ratio} & \textbf{Thuật toán} & \textbf{Robust Accuracy} & \textbf{ASR} & \textbf{Sparsity} & \textbf{PSNR (dB)} \\
        \midrule
        & Sparse-PGD & 56.30\% & 33.14\% & 17.06\% & 32.33 \\
        \textbf{0.1} & GreedyFool & 69.59\% & 17.46\% & 84.99\% & 39.54 \\
        & \cellcolor{yellow!20}\textbf{Proposed-TopkPGD} & \cellcolor{yellow!20}\textbf{72.70\%} & \cellcolor{yellow!20}\textbf{13.66\%} & \cellcolor{yellow!20}\textbf{89.80\%} & \cellcolor{yellow!20}\textbf{40.47} \\
        \midrule
        & Sparse-PGD & 55.00\% & 34.68\% & 9.98\%  & 32.06 \\
        \textbf{0.2} & GreedyFool & 61.70\% & 26.72\% & 70.86\% & 36.59 \\
        & \cellcolor{yellow!20}\textbf{Proposed-TopkPGD} & \cellcolor{yellow!20}\textbf{65.00\%} & \cellcolor{yellow!20}\textbf{22.80\%} & \cellcolor{yellow!20}\textbf{79.48\%} & \cellcolor{yellow!20}\textbf{37.44} \\
        \midrule
        & Sparse-PGD & 56.90\% & 32.42\% & 9.11\%  & 32.13 \\
        \textbf{0.3} & GreedyFool & 56.30\% & 33.14\% & 57.74\% & 34.90 \\
        & \cellcolor{yellow!20}\textbf{Proposed-TopkPGD} & \cellcolor{yellow!20}\textbf{60.30\%} & \cellcolor{yellow!20}\textbf{28.38\%} & \cellcolor{yellow!20}\textbf{69.15\%} & \cellcolor{yellow!20}\textbf{35.67} \\
        \bottomrule
    \end{tabular}%
    }
    \vspace{0.1cm}
    \begin{itemize}
        \item \textbf{Giữ vững độ thưa trước phòng thủ:} Kể cả khi đối mặt với Robust Model, \textbf{Proposed-TopkPGD} vẫn vượt trội hoàn toàn về mặt tàng hình. Tại $k=0.2$, nó đạt độ thưa \textbf{79.48\%} và PSNR \textbf{37.44 dB} trong khi đòn Sparse-PGD sụp đổ xuống còn $9.98\%$.
    \end{itemize}
\end{frame}

% ===============================
% PHẦN 3: THỰC NGHIỆM PHÒNG THỦ
% ===============================
\section{Thực nghiệm Phòng thủ}

% Slide 3.1: Hiệu năng phòng thủ trên Mô hình Chuẩn
\begin{frame}{3.1 Đánh giá Phòng thủ Tiền xử lý (Mô hình Chuẩn)}
    \centering
    \small \textbf{DEFENSE PERFORMANCE ON STANDARD RESNET-18}
    \vspace{0.1cm}
    
    \resizebox{\textwidth}{!}{%
    \begin{tabular}{lcccccc}
        \toprule
        \textbf{Cơ chế phòng thủ} & \textbf{Clean} & \textbf{PGD} & \textbf{Proposed-TopkPGD} & \textbf{Sparse-PGD} & \textbf{GreedyFool} \\
        \midrule
        No Defense & 92.19\% & 0.00\% & 17.19\% & 6.25\% & 4.69\% \\
        Median Filter ($3 \times 3$) & 71.88\% & 43.75\% & \cellcolor{yellow!20}\textbf{68.75\%} & 68.75\% & 57.81\% \\
        Bit Reduction (3-bit) & 82.81\% & 1.56\% & \textbf{50.00\%} & 35.94\% & 32.81\% \\
        JPEG Comp. (Q75) & 84.38\% & 42.19\% & \textbf{60.94\%} & 64.06\% & 56.25\% \\
        Random Noise ($\sigma=0.02$) & 89.06\% & 0.00\% & \textbf{26.56\%} & 6.25\% & 10.94\% \\
        Randomized Sm. & 20.31\% & 15.63\% & \textbf{18.75\%} & 20.31\% & 18.75\% \\
        \bottomrule
    \end{tabular}%
    }
    \vspace{0.2cm}
    \begin{itemize}
        \item \textbf{Hiệu quả của tiền xử lý chống tấn công thưa:} \textbf{Median Filter} và \textbf{JPEG Compression} là "khắc tinh" của các đòn tấn công thưa.
        \item Tuy nhiên, \textbf{Proposed-TopkPGD} thể hiện sức chống chịu (robustness) tốt hơn với các tiền xử lý như Bit Reduction hay Random Noise khi so sánh với Sparse-PGD và GreedyFool.
    \end{itemize}
\end{frame}

% Slide 3.2: Phân tích chuyên sâu khả năng Xuyên Giáp
\begin{frame}{3.2 Phân tích chuyên sâu: Khả năng "Xuyên Giáp" trước Phòng thủ}
    Mặc dù \textbf{Proposed-TopkPGD} có ASR ban đầu thấp hơn SOTA, nhưng nhiễu của nó lại cực kỳ \textbf{khó bị hóa giải (robust)} trước các phép biến đổi ảnh:
    \vspace{0.2cm}
    \begin{itemize}
        \item \textbf{Chống lại Randomized Smoothing:} ASR của GreedyFool giảm thê thảm từ $94.9\% \rightarrow 77.9\%$. Proposed-TopkPGD hầu như không suy suyển ($81.3\% \rightarrow 79.6\%$), chính thức \textbf{vượt mặt SOTA} sau khi qua màng lọc.
        \item \textbf{Chống lại Nén ảnh (JPEG Q75):} GreedyFool mất tới $\sim 64.4\%$ sát thương ($94.9\% \rightarrow 30.5\%$). Thuật toán đề xuất chỉ mất $\sim 50.8\%$, giữ được mức sát thương \textbf{ngang ngửa} GreedyFool dù xuất phát điểm thấp hơn.
        \item \textbf{Cộng hưởng với Random Noise:} Trong khi GreedyFool bị cản phá (ASR giảm), Proposed-TopkPGD lại được "trợ lực" bởi nhiễu ngẫu nhiên, đẩy ASR lên mức \textbf{tuyệt đối 100\%}.
    \end{itemize}
    \begin{block}{Thông điệp cốt lõi (Take-away Message)}
        \textit{"Các phương pháp SOTA cố nhồi nhét cường độ lớn vào số ít pixel nên rất dễ vỡ (fragile) trước bộ lọc. Proposed-TopkPGD phân tán thông minh hơn, không chỉ tàng hình tốt (SSIM cực cao) mà còn sinh ra cấu trúc nhiễu bền vững, kiên cường hơn trong môi trường thực tế."}
    \end{block}
\end{frame}

% Slide 3.3: Hiệu năng phòng thủ trên Mô hình Đối kháng (AT)
\begin{frame}{3.3 Đánh giá Phòng thủ Tiền xử lý (Mô hình AT)}
    \centering
    \small \textbf{DEFENSE PERFORMANCE ON ROBUST RESNET-18 (AT)}
    \vspace{0.1cm}
    
    \resizebox{\textwidth}{!}{%
    \begin{tabular}{lcccccc}
        \toprule
        \textbf{Cơ chế phòng thủ} & \textbf{Clean} & \textbf{PGD} & \textbf{Proposed-TopkPGD} & \textbf{Sparse-PGD} & \textbf{GreedyFool} \\
        \midrule
        No Defense & 84.38\% & 82.81\% & 82.81\% & 84.38\% & 81.25\% \\
        Median Filter ($3 \times 3$) & 75.00\% & 75.00\% & 75.00\% & 75.00\% & 75.00\% \\
        Bit Reduction (3-bit) & 84.38\% & 81.25\% & 82.81\% & 82.81\% & 84.38\% \\
        JPEG Comp. (Q75) & 81.25\% & 81.25\% & 79.69\% & 81.25\% & 79.69\% \\
        Random Noise ($\sigma=0.02$) & 82.81\% & 81.25\% & 82.81\% & 82.81\% & 81.25\% \\
        \bottomrule
    \end{tabular}%
    }
    \vspace{0.2cm}
    \begin{itemize}
        \item \textbf{Sức mạnh áp đảo của Huấn luyện đối kháng (AT):} Bản thân mô hình Robust AT khi không có tiền xử lý đã đạt độ bền vững vô cùng ấn tượng (duy trì $>81\%$ Accuracy dưới hầu hết các loại đòn tấn công thưa).
        \item Các bộ lọc tiền xử lý lúc này không còn mang lại lợi ích lớn, ngược lại còn làm giảm phần nào độ chính xác trên ảnh sạch (Clean Accuracy).
    \end{itemize}
\end{frame}

% Slide 3.4: Hiệu năng phòng thủ trên Mô hình TRADES
\begin{frame}{3.4 Đánh giá Phòng thủ Tiền xử lý (Mô hình TRADES)}
    \centering
    \small \textbf{DEFENSE PERFORMANCE ON TRADES ROBUST MODEL}
    \vspace{0.1cm}
    
    \resizebox{\textwidth}{!}{%
    \begin{tabular}{lcccccc}
        \toprule
        \textbf{Cơ chế phòng thủ} & \textbf{Clean} & \textbf{PGD} & \textbf{Proposed-TopkPGD} & \textbf{Sparse-PGD} & \textbf{GreedyFool} \\
        \midrule
        No Defense & 79.69\% & 51.56\% & \cellcolor{yellow!20}\textbf{71.88\%} & 60.94\% & 70.31\% \\
        Median Filter ($3 \times 3$) & 73.44\% & 56.25\% & \textbf{70.31\%} & 60.94\% & 70.31\% \\
        Bit Reduction (3-bit) & 79.69\% & 53.12\% & \textbf{71.88\%} & 64.06\% & 70.31\% \\
        JPEG Comp. (Q75) & 78.12\% & 54.69\% & \textbf{75.00\%} & 62.50\% & 73.44\% \\
        Random Noise ($\sigma=0.02$) & 79.69\% & 53.12\% & \textbf{73.44\%} & 62.50\% & 70.31\% \\
        \bottomrule
    \end{tabular}%
    }
    \vspace{0.2cm}
    \begin{itemize}
        \item \textbf{Sự kiên cường của Proposed-TopkPGD:} Trên kiến trúc TRADES, phương pháp đề xuất chứng minh được đặc tính **khó bị hóa giải** của nó. Trong khi đòn Sparse-PGD bị cản phá mạnh mẽ, đòn Proposed-TopkPGD vẫn duy trì mức xuyên thấu ổn định.
    \end{itemize}
\end{frame}

% Slide 3.4: Khảo sát trên Mô hình GG-SAT Robust
\begin{frame}{3.4 Khảo sát khả năng phòng thủ của kiến trúc GG-SAT}
    \centering
    \small \textbf{GG-SAT ROBUST MODEL EVALUATION REPORT (1000 samples)}
    \vspace{0.1cm}

    \resizebox{\textwidth}{!}{%
    \begin{tabular}{lccc}
        \toprule
        \textbf{Phương thức tấn công} & \textbf{Tham số cấu hình} & \textbf{Robust Accuracy} & \textbf{ASR} \\
        \midrule
        Clean (No Attack) & - & 88.30\% & 0.00\% \\
        FGSM              & $\epsilon = 8/255$ & 48.70\% & 39.60\% \\
        PGD-10            & $\epsilon = 8/255$ & 36.20\% & 52.10\% \\
        \midrule
        Sparse PGD        & $k = 0.1$ & 73.80\% & 14.50\% \\
        Sparse PGD        & $k = 0.2$ & 63.90\% & 24.40\% \\
        Sparse PGD        & $k = 0.3$ & 57.60\% & 30.70\% \\
        Sparse PGD        & $k = 0.5$ & 48.40\% & 39.90\% \\
        \bottomrule
    \end{tabular}%
    }
    \vspace{0.2cm}
    \begin{itemize}
        \item Mô hình kiến trúc đồ thị GG-SAT có khả năng kháng cự nhiễu khá ấn tượng, vượt mặt mô hình Chuẩn một khoảng xa (Giữ $73.8\%$ Accuracy tại $k=0.1$).
        \item Tuy nhiên, khi đối mặt với nhiễu dày đặc hơn ($k \ge 0.5$), lớp phòng thủ này cũng bắt đầu sụp đổ giống như các mô hình khác.
    \end{itemize}
\end{frame}

% ===============================
% PHẦN 4: HƯỚNG PHÁT TRIỂN
% ===============================
\section{Hướng phát triển}

% Slide 4.1: Hướng phát triển - GGSAT Defense
\begin{frame}{4.1 Hướng phát triển: Phòng thủ với cấu trúc GGSAT}
    \begin{block}{Dự án: Đào tạo lại mô hình với Nhiễu đối kháng thưa}
        \begin{itemize}
            \item \textbf{Ý tưởng cốt lõi:} Lấy chính các mẫu nhiễu thưa hiệu suất cao sinh ra từ thuật toán \textit{Proposed-TopkPGD} để đưa ngược vào quá trình huấn luyện (Adversarial Training), kết hợp sức mạnh kiến trúc đồ thị logic của GGSAT.
            \item \textbf{Mục tiêu nghiên cứu:} Xây dựng một mạng AI tiên tiến có "kháng thể" miễn nhiễm triệt để với các đòn tấn công thưa diện rộng mà ngay cả Median Filter cũng chưa thể bịt kín 100\% lỗ hổng.
            \item \textbf{Tiến độ thực hiện:} \textit{[Chừa chỗ để chèn biểu đồ / bảng kết quả so sánh GG-SAT Defense vs. Standard AT sau khi thu thập đủ thực nghiệm]}
        \end{itemize}
    \end{block}
\end{frame}

\end{document}
